\newpage
\section*{Working Literature Review}

This working review is intended to sit immediately after the Specific Aims page and serve as the starting evidence base for the later Significance, Innovation, and Approach sections. It is deliberately organized around the four questions identified by the PI and links key claims to explicit PubMed IDs.

\subsection*{1. How many patients in the United States may benefit from this technology?}

The addressable population is likely substantial, but it is not captured by a single clean epidemiology paper because the relevant populations span several overlapping clinical groups: patients with heart failure and dyssynchrony who may be considered for CRT or conduction-system pacing, patients with structural heart disease and ventricular tachycardia who may undergo invasive mapping or ablation, and patients with inherited arrhythmia syndromes in whom mechanism and risk stratification remain uncertain. Existing heart-failure epidemiology suggests that millions of adults in the United States are living with heart failure (PMID 29173412; PMID 22227365), while CRT remains underutilized even among eligible patients (PMID 24948569). Sudden cardiac death remains a major public health burden in the United States (PMID 25813838), and Brugada syndrome is a small but clinically important subgroup because of its association with ventricular arrhythmias and sudden death (PMID 31586529; PMID 31993492).

For grant framing, the most defensible near-term statement may be that this technology is relevant to three clinically meaningful groups rather than one artificially merged denominator: (i) patients with ventricular tachycardia being considered for ablation, (ii) heart-failure patients being considered for CRT or conduction-system pacing optimization, and (iii) inherited arrhythmia populations such as Brugada syndrome. A later systematic epidemiology table should estimate each subgroup separately, then provide a justified aggregate range.

\paragraph{Key PMIDs to expand}
PMID 29173412, PMID 22227365, PMID 24948569, PMID 25813838, PMID 31586529, PMID 31993492.

\subsection*{2. What has been done in ECGi, what are its strengths and weaknesses, and where has it been applied?}

Electrocardiographic imaging (ECGi) was developed to reconstruct cardiac electrical information from body-surface recordings and cardiac-torso geometry. Its core strength is that it provides a non-invasive approach to mapping activation from a single beat, which is attractive for ventricular arrhythmia localization and potentially for treatment planning (PMID 34777827). ECGi has also been used to study repolarization abnormalities and non-invasive electrophysiologic phenotypes (PMID 33880931). These strengths make it clinically attractive in settings where invasive mapping is expensive, risky, or incomplete.

The key limitation is that ECGi is fundamentally an inverse problem, and inverse reconstruction from remote body-surface potentials is highly ill-posed (PMID 15649241). In practice, this means that ECGi is sensitive to modeling assumptions, noise, geometry, and regularization choices, and that different internal activation states may project to similar body-surface patterns. In ventricular arrhythmia applications, human accuracy has been variable, and published review work suggests that ECGi may not be sufficiently accurate for highly discrete radiofrequency ablation targeting even if it is useful for broader localization (PMID 34777827). More recent studies continue to evaluate agreement between ECGi systems and invasive maps, reinforcing the need for careful validation rather than overstatement (PMID 37390586).

Current ECGi applications include ventricular arrhythmia localization, repolarization mapping, exploratory substrate characterization, and support for non-invasive treatment paradigms such as radioablation planning. The key translational gap is that ECGi usually reconstructs activity on cardiac surfaces, but does not robustly infer the underlying three-dimensional substrate or distinguish competing internal mechanisms. That gap is precisely where a forward-calibrated mechanistic framework may add value.

\paragraph{Key PMIDs to expand}
PMID 15649241, PMID 34777827, PMID 33880931, PMID 37390586.

\subsection*{3. What are the current tools used to guide CRT and VT ablation, and what are their strengths and weaknesses?}

For CRT and conduction-system pacing, current guidance tools include standard ECG metrics such as QRS duration and morphology, echocardiographic measures, heart-failure phenotype, scar imaging, and increasingly vectorcardiographic or conduction-system metrics. Contemporary guidelines emphasize patient selection, procedural evaluation, response optimization, and follow-up for CRT and conduction-system pacing (PMID 37283271). The major strength of current CRT guidance is that it is clinically familiar and already embedded in routine care. The major weakness is that response remains variable, with a persistent non-responder population and incomplete mechanistic understanding of why a given patient does or does not resynchronize well (PMID 35715087; PMID 37767743; PMID 30091132). Reviews of CRT evolution and conduction-system pacing highlight both the therapeutic promise and the limitations of current selection logic (PMID 33247913).

For VT ablation, the current toolbox includes surface ECG interpretation, invasive electroanatomic mapping, entrainment mapping, pace mapping, substrate mapping, and increasingly pre-procedure or integrated imaging such as late gadolinium enhancement cardiac MRI. Expert consensus documents and reviews support electroanatomic mapping and ablation as core tools in contemporary VT management (PMID 31075787; PMID 31706471). Entrainment mapping remains mechanistically powerful, but is limited by the ability to sustain VT and by hemodynamic tolerance (PMID 28167086). Imaging and scar characterization are promising adjuncts because they may identify conducting channels and slow-conduction regions non-invasively, but they are not yet complete replacements for invasive mapping (PMID 38262674; PMID 36607130).

Taken together, current tools are clinically useful but incomplete. CRT guidance lacks reliable mechanistic predictors of response, and VT ablation still depends heavily on invasive mapping, operator expertise, and indirect interpretation of substrate. This creates a strong rationale for a calibrated forward framework that can integrate imaging and electrical data into a patient-specific mechanism model.

\paragraph{Key PMIDs to expand}
PMID 37283271, PMID 35715087, PMID 37767743, PMID 30091132, PMID 33247913, PMID 31075787, PMID 31706471, PMID 28167086, PMID 38262674, PMID 36607130.

\subsection*{4. What work has been done on molecular markers for CRT and VT ablation success?}

For CRT, the biomarker literature is more mature than for VT ablation. Reviews of CRT biomarkers suggest ongoing interest in fibrosis markers, inflammatory markers, remodeling markers, and electrical phenotypes as predictors of response (PMID 31681965). Specific studies suggest that lower fibrosis biomarker levels may predict better CRT response (PMID 30988339), while newer work is evaluating additional systemic fibrosis-related indices (PMID 39101582). The overall strength of this literature is that it recognizes biological heterogeneity in CRT response; the weakness is that many biomarkers remain inconsistently validated, cohort-specific, or difficult to integrate into routine decision-making without better mechanistic context.

For VT ablation, the evidence base is currently stronger for structural and imaging-derived substrate markers than for truly established circulating molecular markers. Studies increasingly point to scar architecture, conducting channels, lipomatous metaplasia, and inflammatory activity as predictors of arrhythmogenicity or recurrence after ablation (PMID 36607130; PMID 38703163; PMID 33004129; PMID 37819047). This is important for the grant: the VT field appears to have a real translational gap in robust, clinically deployed molecular markers, which strengthens the case for a patient-specific integrative framework rather than a purely biomarker-driven strategy.

The review therefore suggests a nuanced framing. For CRT, biomarker and molecular-predictor work exists but is not yet sufficient for robust individualized guidance. For VT ablation, the field is still dominated by structural and electrophysiologic substrate characterization, with molecular predictors much less mature. This is not a weakness of the proposal; it is part of the unmet need.

\paragraph{Key PMIDs to expand}
PMID 31681965, PMID 30988339, PMID 39101582, PMID 36607130, PMID 38703163, PMID 33004129, PMID 37819047.

\subsection*{Immediate expansion priorities}

The next literature-building step should be to convert each subsection above into a true structured review with:
\begin{itemize}
    \item a clearer search strategy and date limits,
    \item a table of high-priority PMIDs,
    \item one paragraph on strengths of the current field,
    \item one paragraph on key weaknesses and unmet needs,
    \item an explicit bridge sentence linking the literature gap to the proposed ffECG platform.
\end{itemize}
